
\documentclass[italian,bachelor]{unibg}

\title{Integrazione di eBPF in   

 ambiente Android}
\subtitle{Sfide e opportunità per l'ispezione di sistema}
\advisor{Prof. Matthew Rossi}
% \coadvisor{Chiar.mo Prof.~Ed Smith}

\department{Ingegneria Gestionale, dell'Informazione e della Produzione}
\course{Ingegneria Informatica}
\class{L-8}

\author{Mattia Ferrari}
\studentid{1083721}
\academicyear{2024/2025}

\addbibresource{bibliography.bib}
\usepackage{booktabs}
\usepackage{subcaption}
\usepackage{url}
\usepackage{float}
\urlstyle{tt}

\begin{document}
\maketitle
\emptypage

\begin{abstract}
Questa tesi affronta il problema dell'utilizzo di eBPF e bpftrace per il monitoraggio di applicazioni Android a livello kernel. Android, pur basandosi sul kernel Linux, non rende disponibili nativamente gli strumenti di tracing tipici dei sistemi Linux, rendendo necessari adattamenti specifici per poterli utilizzare.

La soluzione proposta si basa sull'esecuzione di bpftrace all'interno di una distribuzione Debian installata nell'emulatore Android Cuttlefish, in esecuzione su Ubuntu. Questo ambiente consente di osservare gli eventi del kernel Android tramite eBPF senza modificare il sistema operativo e senza richiedere hardware fisico.

Il sistema sviluppato comprende sette probe bpftrace parallele per il monitoraggio di syscall, transazioni IPC Binder, ciclo di vita dei processi e traffico di rete, coordinate da un modulo di orchestrazione in Python e analizzate da un modulo post-sessione.

Le sperimentazioni su tre scenari, la navigazione web con Chrome, l'uso della galleria e della fotocamera, e l'attivazione della modalità aereo, mostrano che il monitoraggio parallelo delle probe permette di ricostruire il comportamento interno del sistema operativo in modo non ottenibile da una singola fonte di dati. I risultati evidenziano l'architettura di isolamento del renderer Chromium, la firma Binder specifica dell'accesso all'hardware fotografico e la sequenza di disattivazione e riconnessione dello stack WiFi.
\end{abstract}

\emptypage
\toc
\emptypage

\clearpage
\pagenumbering{arabic}

\chapter{Introduzione}
\section{Scenario applicativo}
Negli ultimi anni i dispositivi mobili hanno assunto un ruolo centrale nell’ambito dell’informatica personale e professionale.
Tra i sistemi operativi per dispositivi mobili, Android rappresenta una delle piattaforme più diffuse a livello mondiale~\cite{android_marketshare} e 
viene utilizzato in una grande varietà di dispositivi, dagli smartphone ai sistemi embedded. La complessità crescente delle 
applicazioni e del sistema operativo rende sempre più importante disporre di strumenti in grado di analizzare e comprendere il 
comportamento delle applicazioni durante l’esecuzione.

Il monitoraggio delle attività delle applicazioni è fondamentale in diversi ambiti, tra cui il debugging del software, 
l’analisi delle prestazioni e la sicurezza informatica. In particolare, la possibilità di osservare le interazioni tra le 
applicazioni e il sistema operativo consente di individuare anomalie di funzionamento, colli di bottiglia prestazionali e 
potenziali comportamenti malevoli.  

Tradizionalmente il tracing del sistema operativo viene realizzato mediante strumenti specifici che permettono di raccogliere 
informazioni sull’esecuzione dei processi e sulle attività del kernel. Tali strumenti possono tuttavia introdurre un significativo 
overhead computazionale oppure richiedere modifiche al sistema operativo, rendendone difficile l’utilizzo in ambienti reali. 
Inoltre, molti strumenti di analisi disponibili nei sistemi operativi tradizionali, e in particolare in ambiente Linux, non 
sono direttamente utilizzabili su Android senza adattamenti specifici.

Nel 2014 è stato introdotto l’extended Berkeley Packet Filter (eBPF) in Linux~\cite{ebpf_kernel}, una tecnologia che 
permette di eseguire piccoli programmi in modo sicuro ed efficiente all’interno del sistema operativo. Questo consente di osservare 
il comportamento delle applicazioni e del sistema senza la necessità di modifiche al codice sorgente o dello sviluppo di moduli dedicati. 
Per queste ragioni eBPF è diventato uno strumento sempre più utilizzato per attività di tracing e osservabilità nei sistemi Linux~\cite{ebpf_observability}.
\begin{figure}[h]         
    \centering
    \includegraphics[width=0.8\textwidth]{images/architettura-ebpf.png}
    \caption{Architettura di eBPF nel kernel Linux~\cite{ebpf_observability}}
    \label{fig:architettura-ebpf}
\end{figure}

\section{Criticità}
Tra gli strumenti che permettono di utilizzare eBPF in modo pratico, bpftrace costituisce una soluzione particolarmente efficace per 
il tracing delle applicazioni e del sistema operativo.\texttt{ bpftrace}~\cite{bpftrace_docs} mette a disposizione un linguaggio di scripting ad alto livello che 
consente di definire sonde e azioni associate agli eventi osservati, permettendo di raccogliere informazioni durante l’esecuzione senza 
sviluppare programmi eBPF completi. Questo approccio semplifica notevolmente la realizzazione di strumenti di analisi dinamica.

Tuttavia, l’utilizzo di \texttt{ bpftrace} in ambiente Android presenta diverse difficoltà pratiche. In molti dispositivi l’accesso alle 
funzionalità di basso livello del sistema è limitato e spesso non sono disponibili gli strumenti e le librerie necessari per l’esecuzione di programmi basati su eBPF. Inoltre, l’installazione diretta di tali strumenti sul sistema Android può risultare complessa 
o non praticabile, soprattutto al di fuori di ambienti di sviluppo, quando non si ha accesso a certi privilegi che consentono 
di effettuare l'installazione e l'uso di tali strumenti. Ottenere tali 
permessi richiederebbe di effettuare il rooting del dispositivo o di sostituire la release 
Android installata, operazioni spesso non praticabili in contesti reali.

\begin{table}[h]
    \centering
    \begin{tabular}{lcc}
        \toprule
        & \textbf{Emulatore} & \textbf{Dispositivo di Produzione} \\
        \midrule
        \textbf{Accesso root}           & Disponibile     & Richiede rooting   \\
        \textbf{Installazione strumenti} & Possibile       & Richiede rooting  \\
        \textbf{Modifica della ROM}     & Non necessaria  & Spesso necessaria \\
        \textbf{Utilizzo di bpftrace}   & Fattibile       & Non praticabile   \\
        \textbf{Riproducibilità}        & Alta            & Bassa             \\
        \bottomrule
    \end{tabular}
    \caption{Confronto tra ambiente di test e dispositivo di produzione Android.}
    \label{fig:confronto-ambienti}
\end{table}

Queste limitazioni rendono difficile l’impiego delle tecnologie basate su eBPF per l’analisi del comportamento delle applicazioni Android, 
nonostante il loro potenziale per il monitoring del sistema. Risulta quindi necessario definire un ambiente di lavoro che consenta l’utilizzo 
di bpftrace in modo stabile e riproducibile, permettendo allo stesso tempo l’osservazione delle attività delle applicazioni durante l’esecuzione.

Il problema affrontato in questa tesi consiste quindi nel rendere possibile l’utilizzo di bpftrace per il tracing delle applicazioni 
in ambiente Android, individuando una configurazione che consenta di superare le limitazioni tipiche del sistema e di ottenere dati significativi sull’esecuzione delle applicazioni.

\section{Approccio proposto}
Per superare le difficoltà legate all’utilizzo degli strumenti basati su eBPF in ambiente Android, in questo lavoro è stata 
sviluppata una soluzione che permette di eseguire lo strumento bpftrace in un ambiente controllato e riproducibile, mantenendo 
allo stesso tempo la possibilità di osservare il comportamento delle applicazioni Android.

La soluzione proposta si basa sulla creazione di un ambiente di sviluppo virtualizzato in cui sia possibile eseguire Android e 
utilizzare strumenti tipicamente disponibili nei sistemi Linux tradizionali. A questo scopo è stato utilizzato l’emulatore 
Android Cuttlefish~\cite{cuttlefish_docs}, eseguito su sistema operativo Ubuntu, che consente di disporre di un ambiente Android completo e 
configurabile per attività di sviluppo e sperimentazione.

All’interno dell’ambiente Android è stata installata una distribuzione Linux minimale basata su Debian~\cite{debian_distrib}. Questa distribuzione è 
stata utilizzata come ambiente di lavoro per l’esecuzione di bpftrace e degli strumenti necessari al tracing. L’utilizzo di una 
distribuzione Linux separata ha permesso di evitare la compilazione nativa di bpftrace su Android e di semplificare l’
installazione delle dipendenze necessarie.

In questo modo è stato possibile realizzare un ambiente che combina la flessibilità degli strumenti Linux con la possibilità di 
osservare il comportamento di applicazioni eseguite su Android. Su questa base è stato sviluppato un insieme di script e 
strumenti software che permettono di eseguire sessioni di tracing e di raccogliere automaticamente i dati prodotti da bpftrace.
\begin{figure}[h]
      \centering
      \includegraphics[width=0.6\textwidth]{images/architettura-soluzione.png}
      \caption{Architettura della soluzione proposta: Ubuntu, Cuttlefish e ambiente Debian.}
      \label{fig:architettura-soluzione}
\end{figure}


\chapter{Tecnologie di base}
\section{Linux}
Linux è una famiglia di sistemi operativi open-source Unix-like, basati sull'omonimo kernel, ampiamente utilizzato sia in ambito accademico sia industriale~\cite{linux_kernel} grazie alla sua stabilità, 
flessibilità e natura open source. Il kernel gestisce le risorse hardware e fornisce ai programmi 
un’interfaccia standardizzata attraverso le system call, mentre l’insieme dei programmi in spazio utente costituisce l’ambiente operativo vero e proprio. 
Ciò che rende Linux particolarmente adatto allo sviluppo software e alle attività di sperimentazione è la sua natura open source e free: il codice sorgente è liberamente accessibile, modificabile e distribuibile, il che consente un controllo dettagliato del sistema e facilita l'accesso a un ecosistema ampio di strumenti di analisi e debugging.

\begin{figure}[h]
    \centering
    \includegraphics[width=0.8\textwidth]{images/linux-architettura.png}
    \caption{Architettura del sistema Linux: kernel space, user space e system call~\cite{linux_architettura_img}.}
    \label{fig:linux-architettura}
\end{figure}

In questo lavoro è stato utilizzato Ubuntu 22.04 LTS, una distribuzione Linux ampiamente diffusa. Ubuntu è stato scelto per la disponibilità di pacchetti software aggiornati e per la semplicità di configurazione dell’ambiente di 
sviluppo. Il sistema Ubuntu ha costituito la piattaforma principale su cui sono stati installati gli strumenti necessari allo svolgimento della tesi, 
tra cui l’ambiente di virtualizzazione Android Cuttlefish e i componenti utilizzati per il tracing basato su eBPF. L’utilizzo di una distribuzione Linux 
standard ha permesso di disporre di un ambiente stabile e riproducibile, facilitando l’installazione delle dipendenze e la gestione degli strumenti software impiegati nelle sperimentazioni.

\section{Android Cuttlefish}
Android~\cite{android_aosp} è un sistema operativo per dispositivi mobili basato sul kernel Linux e progettato per supportare l’esecuzione di applicazioni in un ambiente isolato e controllato. Il sistema è organizzato in una struttura a livelli che comprende il kernel Linux, le librerie di sistema, il runtime Android, il framework delle applicazioni e il livello delle applicazioni utente. Il kernel Linux costituisce il livello più basso dell’architettura e si occupa della gestione delle risorse hardware, della memoria, dei processi, dei dispositivi di input/output e delle comunicazioni di rete. I livelli superiori forniscono invece i servizi necessari all’esecuzione delle applicazioni e definiscono l’ambiente software tipico della piattaforma Android.

Ogni applicazione Android viene eseguita in uno spazio isolato rispetto alle altre applicazioni, generalmente associato a un utente distinto. Questo meccanismo di isolamento contribuisce alla sicurezza del sistema ma limita l’accesso diretto alle risorse e alle funzionalità di basso livello. Inoltre, molte componenti tipiche delle distribuzioni Linux tradizionali non sono presenti oppure risultano fortemente ridotte, rendendo più complesso l’utilizzo di strumenti di analisi progettati per ambienti Linux standard.

Nonostante Android utilizzi il kernel Linux, l’ambiente software differisce in modo significativo da quello di una distribuzione Linux tradizionale. In particolare, l’organizzazione del file system, la disponibilità delle librerie e la gestione dei permessi rendono spesso necessario un adattamento degli strumenti sviluppati per sistemi Linux convenzionali. Questo aspetto è particolarmente rilevante nel caso degli strumenti di tracing, che richiedono generalmente l’accesso a funzionalità di basso livello del sistema operativo.

\begin{figure}[h]
    \centering
    \includegraphics[width=0.7\textwidth]{images/android-architettura.png}
    \caption{Architettura a livelli del sistema operativo Android~\cite{android_architettura_img}.}
    \label{fig:android-architettura}
\end{figure}

Per lo sviluppo e la sperimentazione è stato utilizzato Android Cuttlefish~\cite{cuttlefish_docs}, un emulatore ufficiale della piattaforma Android progettato per l’esecuzione su sistemi Linux. Cuttlefish esegue il sistema Android all’interno di una macchina virtuale che riproduce i principali componenti di un dispositivo reale, inclusi kernel, sistema operativo e servizi applicativi, senza la necessità di hardware fisico.

In questo progetto l’emulatore è stato eseguito su Ubuntu, che ha costituito la piattaforma principale per l’installazione e la configurazione degli strumenti di tracing. L’ambiente virtualizzato ha permesso di controllare la configurazione del sistema e di ripetere le sperimentazioni in modo sistematico.
\section{Tracing, eBPF e bpftrace}
\subsection{Tracing delle applicazioni}
Il tracing delle applicazioni è una tecnica di analisi che consiste nella raccolta sistematica di informazioni durante l’esecuzione di un programma o di un sistema operativo, 
con l’obiettivo di osservare il comportamento dinamico del software e le sue interazioni con il sistema. A differenza dell’analisi statica, che si basa sull’esame del codice 
sorgente o dei file eseguibili, il tracing consente di studiare ciò che avviene realmente durante l’esecuzione, permettendo di ottenere informazioni temporali sugli eventi generati dal sistema.

Nel contesto dei sistemi operativi moderni, il tracing può essere effettuato a diversi livelli di osservazione. A livello applicativo è possibile monitorare le operazioni 
effettuate da un processo, come la creazione di nuovi processi, l’accesso ai file o le comunicazioni di rete. A livello di sistema operativo è invece possibile osservare 
eventi gestiti dal kernel, come le chiamate di sistema, la schedulazione dei processi o le operazioni di input/output. Questa possibilità di osservare eventi a diversi 
livelli rende il tracing uno strumento fondamentale per comprendere il funzionamento complessivo di un sistema software.

Le tecniche di tracing vengono utilizzate principalmente per tre scopi. Il primo è il debugging del software, dove l’osservazione dettagliata delle operazioni eseguite 
da un programma permette di individuare errori o comportamenti inattesi. Il secondo è l’analisi delle prestazioni, in cui il tracing consente di identificare colli di 
bottiglia e di comprendere come le risorse del sistema vengono utilizzate dalle applicazioni. Il terzo ambito è quello della sicurezza informatica, dove l’osservazione 
delle attività dei processi permette di individuare comportamenti anomali o potenzialmente malevoli.

Tradizionalmente il tracing in ambiente Linux è stato realizzato tramite strumenti basati su meccanismi come \texttt{ptrace}~\cite{ptrace_man}, utilizzati ad esempio da programmi come \texttt{strace}~\cite{strace_docs}, oppure 
tramite infrastrutture di tracing integrate nel kernel come \texttt{ftrace}~\cite{ftrace_docs} e \texttt{perf}~\cite{perf_docs}. Questi strumenti permettono di ottenere informazioni dettagliate sul comportamento del sistema, 
ma possono risultare complessi da utilizzare oppure introdurre un overhead significativo durante l’esecuzione.

\subsection{Tracing mediante eBPF}
Una delle tecnologie più recenti per il tracing nei sistemi Linux è rappresentata dall’extended Berkeley Packet Filter (eBPF)~\cite{ebpf_kernel}. eBPF permette di eseguire piccoli programmi 
all’interno del sistema operativo in modo controllato e sicuro, consentendo di intercettare eventi durante l’esecuzione senza modificare il codice del kernel o delle applicazioni.

Il funzionamento di eBPF si basa sull’inserimento dinamico di programmi che vengono eseguiti quando si verificano determinati eventi, come l’ingresso in una funzione del kernel o 
l’attivazione di un tracepoint. Prima di essere caricati, i programmi eBPF vengono verificati da un componente del sistema operativo chiamato verifier, che ne controlla la sicurezza 
e garantisce che non possano compromettere la stabilità del sistema. Questo meccanismo permette di eseguire codice in modo sicuro anche all’interno di parti critiche del sistema operativo.

\begin{figure}[h]
    \centering
    \includegraphics[width=0.7\textwidth]{images/ebpf-diagram.png}
    \caption{Flusso di caricamento ed esecuzione di un programma eBPF nel kernel Linux~\cite{ebpf_observability}.}
    \label{fig:ebpf-flusso}
\end{figure}

Rispetto alle tecniche di tracing tradizionali, eBPF presenta alcune caratteristiche distintive. In primo luogo consente di raccogliere informazioni con un overhead generalmente 
ridotto, poiché i programmi vengono eseguiti direttamente nel contesto dell’evento osservato. Inoltre, le sonde possono essere attivate, disattivate e ridefinite dinamicamente senza riavviare o modificare il sistema operativo, il che permette di adattare l'osservazione alle proprie esigenze senza alcun impatto sul sistema monitorato.

Grazie a queste caratteristiche eBPF rappresenta oggi una tecnologia particolarmente adatta per il tracing dinamico dei sistemi Linux e costituisce la base di numerosi strumenti moderni 
di osservabilità.

\subsection{bpftrace}

Per utilizzare eBPF in modo pratico sono stati sviluppati diversi strumenti software. Tra questi, \texttt{bpftrace}~\cite{bpftrace_docs} rappresenta uno degli strumenti più diffusi per la realizzazione di script di tracing basati su eBPF.

\texttt{bpftrace} è un linguaggio di scripting e un interprete che permette di definire programmi di tracing in forma compatta e leggibile. Gli script vengono tradotti automaticamente in programmi eBPF ed eseguiti dal sistema operativo, semplificando notevolmente lo sviluppo rispetto alla scrittura diretta di codice eBPF.

Un programma \texttt{bpftrace} è costituito da una o più probe, cioè punti di osservazione associati a eventi specifici. Ogni probe è composta da due parti principali:

\begin{itemize}
\item la definizione dell’evento da osservare
\item le azioni da eseguire quando l’evento si verifica
\end{itemize}

La struttura generale di una probe è illustrata in Figura~\ref{fig:bpftrace-probe}, che evidenzia la separazione tra la definizione dell'evento da osservare e il blocco delle azioni da eseguire.

\begin{figure}[h]
    \centering
    \includegraphics[width=0.9\textwidth]{images/bpftrace-probe.png}
    \caption{Struttura di una probe bpftrace: definizione dell'evento e blocco delle azioni.}
    \label{fig:bpftrace-probe}
\end{figure}

\texttt{bpftrace} supporta diversi tipi di probe, che permettono di osservare eventi a diversi livelli del sistema. Nel presente lavoro sono stati utilizzati esclusivamente i tracepoint, in quanto rappresentano la soluzione più stabile e compatibile tra diverse versioni del kernel.

I tracepoint sono punti di osservazione statici definiti all'interno del sistema operativo. A differenza di altre tipologie di probe, la loro interfaccia è stabile tra diverse versioni del kernel, il che li rende particolarmente adatti a un ambiente come Android, dove la versione del kernel può variare. Ogni tracepoint espone un insieme strutturato di argomenti che descrivono l'evento osservato, consentendo di accedere direttamente a informazioni quali il processo coinvolto, i parametri della chiamata e il risultato dell'operazione. I tracepoint disponibili nel sistema utilizzato durante le sperimentazioni sono illustrati in Figura~\ref{fig:tracepoint-list}.

\begin{figure}[h]
    \centering
    \includegraphics[width=0.6\textwidth]{images/tracepoint-list.png}
    \caption{Lista dei tracepoint presenti nell'ambiente Android Cuttlefish e utilizzati durante le sperimentazioni.}
    \label{fig:tracepoint-list}
\end{figure}

\texttt{bpftrace} supporta anche altri due tipi di probe. Le kprobe consentono di inserire sonde dinamiche all'ingresso o all'uscita di funzioni del kernel, offrendo grande flessibilità anche in punti privi di tracepoint predefiniti; la loro stabilità è tuttavia inferiore, poiché dipendono dai nomi interni delle funzioni del kernel che possono variare tra versioni diverse. Le uprobe operano invece in spazio utente e permettono di tracciare funzioni di librerie o applicazioni specifiche senza modificarne il codice, risultando utili quando l'obiettivo è osservare il comportamento interno di un singolo processo piuttosto che eventi di sistema.

\chapter{Implementazione}
\section{Architettura generale}

Il sistema è composto da quattro componenti principali che collaborano secondo un'architettura a pipeline: le probe \texttt{bpftrace} raccolgono eventi direttamente dal kernel, \texttt{monitor.py} li normalizza e persiste, \texttt{summary.py} li analizza in fase post-sessione, e \texttt{launcher.sh} orchestra l'intera interazione con l'utente. L'architettura complessiva è illustrata in Figura~\ref{fig:architettura-sistema}.

\begin{figure}[h]
    \centering
    \includegraphics[width=0.8\textwidth]{images/architettura-sistema.png}
    \caption{Architettura del sistema: flusso di dati tra probe bpftrace, monitor.py, summary.py e launcher.sh.}
    \label{fig:architettura-sistema}
\end{figure}

Il flusso di dati segue un percorso ben definito: le probe producono eventi strutturati in formato JSON su stdout, \texttt{monitor.py} li consuma riga per riga e li persiste in un file \texttt{.jsonl}, e \texttt{summary.py} legge questo file al termine della sessione per produrre report e statistiche aggregate. \texttt{launcher.sh} funge da punto di ingresso unico che coordina l'avvio delle probe, la durata della sessione e la generazione del report finale.

Questa separazione delle responsabilità risponde a un principio architetturale preciso: ciascun componente è progettato per fare una sola cosa e farlo bene, rimanendo testabile e utilizzabile in modo indipendente. Le interfacce di comunicazione tra i componenti sono deliberatamente semplici — stdout/stdin per il flusso di eventi, file JSON per la configurazione e la persistenza — in modo da minimizzare le dipendenze e facilitare la manutenzione.

\subsection{Le probe bpftrace}

Le probe operano interamente in spazio kernel, con accesso diretto alle strutture dati del kernel quali \texttt{task\_struct}~\cite{linux_task_struct}, e producono output in formato JSON su stdout. La scelta di JSON come formato di output permette a bpftrace di comportarsi come un produttore di eventi strutturati che il livello superiore può consumare senza logica di parsing ad hoc. Ogni evento emesso segue uno schema comune con i campi: \texttt{ts\_ns} (timestamp in nanosecondi), \texttt{ts} (orario leggibile), \texttt{type} (categoria dell'evento), \texttt{event} (nome specifico dell'evento), \texttt{pid}, \texttt{ppid}, \texttt{tid}, \texttt{uid}, \texttt{comm} (nome del processo, massimo 15 caratteri imposto dal kernel) e un oggetto \texttt{data} contenente i campi specifici della probe.

\subsection{Il modulo di orchestrazione}

\texttt{monitor.py} è il cuore del sistema. Si occupa di avviare bpftrace come sottoprocesso, gestire il ciclo di vita della sessione di raccolta dati, normalizzare e validare gli eventi ricevuti, arricchirli dove necessario tramite il filesystem \texttt{/proc} e salvarli su disco in formato JSON Lines (\texttt{.jsonl}). Opera come consumer dello stream stdout di bpftrace, consumando gli eventi riga per riga in modo sincrono, mentre un thread dedicato drena contemporaneamente stderr per separare gli errori diagnostici.

\subsection{Il modulo di analisi}

\texttt{summary.py} è il modulo di analisi post-sessione. Legge il file \texttt{events.jsonl} prodotto da \texttt{monitor.py}, calcola statistiche aggregate su syscall, processi, Binder IPC e attività di rete, e genera un report testuale sia in formato plain text che Markdown. Il modulo produce inoltre un file \texttt{index.json} all'interno della directory di sessione, contenente tutti i dati analitici in forma strutturata.

\subsection{L'interfaccia utente}

\texttt{launcher.sh} è lo script bash che funge da interfaccia utente per il sistema. Offre un menu interattivo a colori che permette di scegliere quale probe eseguire, per quanto tempo, e se generare automaticamente un report al termine della sessione. Supporta anche una modalità non interattiva tramite flag da riga di comando (\texttt{-p} per specificare la probe, \texttt{-t} per il timer).

% -------------------------------------------------------
\section{Considerazioni progettuali}

Alcune scelte progettuali meritano una riflessione esplicita. La separazione tra la probe bpftrace (che produce dati grezzi) e \texttt{monitor.py} (che normalizza, arricchisce e persiste) risponde a un principio di separazione delle responsabilità: bpftrace è ottimizzato per la raccolta ad alta frequenza di eventi in spazio kernel, mentre Python è più adatto per la logica applicativa complessa come la lettura di \texttt{/proc} o la gestione delle sessioni. Un'unica componente che facesse entrambe le cose sarebbe più fragile e difficile da manutenere.

La scelta di correlare \texttt{sys\_enter} e \texttt{sys\_exit} nelle probe di rete (invece di usare solo \texttt{sys\_enter}) aggiunge la dimensione temporale all'osservazione: non solo si sa che una syscall è stata invocata, ma anche quanto ha impiegato e se ha avuto successo. Questa informazione è preziosa per l'analisi delle prestazioni e per il rilevamento di anomalie basato su latenza.

La suddivisione dell'attività di rete in quattro probe separate (invece di una sola probe che monitora tutte le syscall di rete) è una scelta di modularità che permette all'utente di attivare solo il sottoinsieme di informazioni di cui ha bisogno, riducendo il volume di dati prodotti e l'overhead sul sistema monitorato.

Infine, la presenza di \texttt{launcher.sh} come wrapper bash — invece di integrare la gestione del timer direttamente in \texttt{monitor.py} — riflette la filosofia Unix di strumenti semplici e componibili~\cite{unix_philosophy}: \texttt{monitor.py} si occupa solo di monitorare, il launcher si occupa dell'orchestrazione e dell'esperienza utente, e \texttt{summary.py} si occupa solo dell'analisi. Questa separazione rende ciascun componente testabile e utilizzabile indipendentemente.


% -------------------------------------------------------
\section{Le probe}

Il progetto include sette probe bpftrace, ognuna dedicata a osservare una specifica classe di eventi di sistema. La directory \texttt{probes/} contiene tutti gli script \texttt{.bt}, mentre la loro configurazione — metadati, codice identificativo, tipo di evento atteso — è centralizzata nel file \texttt{config/probes\_map.json}, che funge da registro consultato sia da \texttt{monitor.py} che da \texttt{launcher.sh}.

\subsection{process\_lifecycle.bt — Ciclo di vita dei processi}

Questa probe monitora le tre fasi fondamentali del ciclo di vita di un processo: la creazione (\textit{fork}), l'esecuzione di un nuovo programma (\textit{exec}) e la terminazione (\textit{exit}). Per farlo, aggancia tre tracepoint dello scheduler del kernel Linux: \texttt{sched:sched\_process\_fork}, \texttt{sched:sched\_process\_exec} e \texttt{sched:sched\_process\_exit}.

La scelta di questi tracepoint è motivata dalla loro stabilità e dalla ricchezza di informazioni che espongono. I tracepoint \texttt{sched\_process\_*} sono presenti e stabili in tutte le versioni del kernel Linux utilizzate da Android, e forniscono direttamente i metadati del processo coinvolto senza necessità di accedere a strutture kernel aggiuntive. Rispetto all’alternativa di intercettare le syscall \texttt{fork} o \texttt{execve} tramite \texttt{raw\_syscalls}, questi tracepoint vengono attivati in un momento del ciclo di vita del processo in cui le strutture dati del kernel sono già aggiornate e consistenti, riducendo il rischio di osservare stati intermedi.

Un aspetto rilevante è la raccolta del \texttt{ppid} (Process ID del processo padre), ottenuto tramite un cast esplicito alla struttura kernel \texttt{task\_struct}: il campo \texttt{real\_parent->tgid} fornisce l'identificativo del gruppo di thread del processo genitore, informazione fondamentale per ricostruire l'albero genealogico dei processi. Il campo \texttt{comm}, recuperato direttamente dal kernel, è soggetto al limite di 15 caratteri imposto dal kernel Linux; per questa ragione, \texttt{monitor.py} arricchisce gli eventi di tipo \texttt{exec} leggendo il file \texttt{/proc/<pid>/cmdline} prima che il processo termini, ottenendo così la riga di comando completa.

Gli eventi generati da questa probe sono utilizzati da \texttt{summary.py} per costruire l'albero dei processi e correlare l'attività dei diversi componenti del sistema Android.

\begin{figure}[h]
    \centering
    \includegraphics[width=0.6\textwidth]{images/process-lifecycle.png}
    \caption{Diagramma del ciclo di vita di un processo Linux}
    \label{fig:process-lifecycle}
\end{figure}

\subsection{syscalls.bt — Monitoraggio delle system call}

La probe \texttt{syscalls.bt} utilizza il tracepoint generico \texttt{raw\_syscalls:sys\_enter}, che viene attivato all'ingresso di \textit{qualsiasi} system call invocata nel sistema. Per limitare l'osservazione alle sole syscall di interesse, la probe applica un filtro inline sul campo \texttt{args->id}, che contiene il numero identificativo della syscall:

\begin{itemize}
    \item \texttt{execve} (id 59): esecuzione di un nuovo programma;
    \item \texttt{openat} (id 257): apertura di un file;
    \item \texttt{connect} (id 42): tentativo di connessione di rete.
\end{itemize}

Questo approccio — intercettare tutto e filtrare — è preferibile rispetto all'uso di tracepoint specifici per syscall (come \texttt{tracepoint:syscalls:sys\_enter\_execve}) perché consente di gestire tutte le syscall di interesse con un unico handler, riducendo la duplicazione del codice e semplificando la gestione delle mappe temporanee condivise.

La probe implementa una decodifica semantica degli argomenti per ciascuna syscall. Per \texttt{execve} viene letto il percorso del binario da eseguire; per \texttt{openat} il percorso del file da aprire; per \texttt{connect} viene ispezionata la struttura \texttt{sockaddr} per determinare la famiglia di indirizzi (\texttt{AF\_UNIX}, \texttt{AF\_INET} o \texttt{AF\_INET6}) e decodificare l'indirizzo di destinazione.

Un aspetto critico da considerare nella lettura di questi parametri è il problema del \textit{time-of-check to time-of-use} (TOCTOU)~\cite{toctou}. I percorsi dei file letti da \texttt{execve} e \texttt{openat}, così come gli indirizzi letti dalla struttura \texttt{sockaddr} in \texttt{connect}, risiedono nella memoria del processo utente. Nel momento in cui la probe li legge tramite \texttt{uptr()}, il processo potrebbe in teoria modificarli prima che il kernel li utilizzi effettivamente. Questo significa che il valore osservato dalla probe potrebbe non corrispondere al valore effettivamente usato dal kernel per eseguire l'operazione. Per scopi di tracing comportamentale e analisi forense questo margine di incertezza è generalmente accettabile, ma è importante tenerne conto quando si interpretano i dati raccolti, in particolare in contesti di sicurezza dove un attaccante potrebbe sfruttare questa finestra temporale.

Un aspetto tecnico degno di nota riguarda la gestione di \texttt{AF\_INET}: la funzione \texttt{ntop()} di bpftrace, necessaria per convertire un indirizzo IPv4 in formato testuale, non può essere assegnata a una variabile ma può essere utilizzata solo inline in un'istruzione \texttt{printf}. Per questo motivo la probe gestisce il caso \texttt{AF\_INET} in un ramo separato del codice, emettendo l'evento con una \texttt{printf} dedicata che include \texttt{ntop()} direttamente nell'argomento stringa di formato, e ritorna immediatamente con \texttt{return} senza passare dal \texttt{printf} generico.

Per evitare memory leak nelle mappe bpftrace utilizzate come variabili temporanee indicizzate per \texttt{tid}, la probe include un handler del tracepoint \texttt{sched:sched\_process\_exit} che cancella i valori residui, e un handler \texttt{END} che pulisce le mappe al termine.

\subsection{Monitoraggio IPC Binder}

Binder~\cite{binder_ipc} è il meccanismo di Inter-Process Communication (IPC) fondamentale di Android. Ogni chiamata a un servizio di sistema, ogni interazione tra applicazioni e framework Android, ogni chiamata a un Content Provider passa attraverso Binder. Monitorare Binder significa quindi avere visibilità sul tessuto connettivo dell'intero sistema operativo Android.

La scelta di usare i tracepoint del driver Binder è obbligata: non esistono syscall dirette che espongano le transazioni Binder in modo strutturato, poiché tutta la comunicazione avviene tramite il driver \texttt{/dev/binder} attraverso \texttt{ioctl}. I tracepoint \texttt{binder:binder\_transaction\_*} sono l'unico punto di osservazione che fornisce informazioni strutturate sulle transazioni a livello kernel, senza necessità di intercettare e decodificare le \texttt{ioctl} manualmente.

La probe aggancia tre tracepoint esposti dal driver kernel del Binder:

\begin{itemize}
    \item \texttt{binder:binder\_transaction}: generato all'avvio di una transazione, cattura mittente, destinatario e metadati della chiamata;
    \item \texttt{binder:binder\_transaction\_received}: generato quando il destinatario riceve la transazione, permette di correlare invio e ricezione tramite il campo \texttt{debug\_id};
    \item \texttt{binder:binder\_transaction\_alloc\_buf}: generato quando viene allocato il buffer dati, fornisce la dimensione del payload trasferito.
\end{itemize}

Il campo \texttt{flags} viene decodificato per determinare se la transazione è one-way (asincrona, bit \texttt{0x01} dei flag impostato) o sincrona. Il \texttt{ppid} è recuperato tramite cast a \texttt{task\_struct} come nelle altre probe.

La correlazione tra i tre eventi tramite \texttt{debug\_id} è fondamentale: solo incrociando \texttt{binder\_transaction} (che ha i dati del mittente), \texttt{binder\_transaction\_alloc\_buf} (che ha la dimensione del payload) e \texttt{binder\_transaction\_received} (che ha i dati del destinatario) è possibile ricostruire un'immagine completa di ogni transazione IPC. Questa correlazione viene eseguita da \texttt{summary.py} nella fase di analisi.

\begin{figure}[h]
    \centering
    \includegraphics[width=0.8\textwidth]{images/inder-correlazione.png}
    \caption{Schema della correlazione tra i tre tracepoint Binder tramite il campo \texttt{debug\_id}.}
    \label{fig:binder-correlazione}
\end{figure}

\subsection{Probe di rete}

L'attività di rete è monitorata tramite quattro probe specializzate, ciascuna dedicata a una syscall diversa. La scelta di osservare le syscall di rete tramite il tracepoint generico \texttt{raw\_syscalls:sys\_enter/sys\_exit} — filtrato per numero di syscall — invece di tracepoint specifici è motivata dalla necessità di correlare ingresso e uscita della stessa syscall per calcolare la latenza e leggere il valore di ritorno. Questa tecnica di correlazione richiede di salvare il timestamp e il contesto all'ingresso in una mappa indicizzata per \texttt{tid}, e di recuperarli all'uscita. La suddivisione in probe separate è inoltre una scelta progettuale che offre flessibilità: è possibile attivare solo il monitoraggio delle connessioni in uscita, o solo il volume di dati inviati, o solo i servizi in ascolto, senza attivare l'intero set.

Anche per le probe di rete vale la considerazione TOCTOU già discussa per \texttt{syscalls.bt}: gli indirizzi letti dalla struttura \texttt{sockaddr} in memoria utente potrebbero in teoria essere modificati dal processo tra il momento della lettura da parte della probe e il momento in cui il kernel utilizza effettivamente il valore.

\noindent\textbf{net\_send.bt}; traccia la syscall \texttt{sendto} (id 44). Utilizza la tecnica di correlazione \texttt{sys\_enter}/\texttt{sys\_exit}: all'ingresso della syscall vengono salvate in mappe bpftrace indicizzate per \texttt{tid} le informazioni di contesto (timestamp, ppid, file descriptor, dimensione richiesta); all'uscita viene calcolata la latenza come differenza di timestamp in nanosecondi e convertita in microsecondi. Il valore di ritorno \texttt{args->ret} corrisponde al numero di byte effettivamente inviati, che può differire dalla dimensione richiesta in caso di invio parziale.


\noindent\textbf{net\_recv.bt}; funziona in modo simmetrico rispetto a \texttt{net\_send.bt} ma aggancia \texttt{recvfrom} (id 45). Registra la dimensione del buffer allocato per la ricezione e il numero di byte effettivamente ricevuti, oltre alla latenza della syscall.


\noindent\textbf{net\_bind.bt}; traccia \texttt{bind} (id 49), la syscall con cui un processo richiede al kernel di associare un socket a un indirizzo locale, tipicamente per mettersi in ascolto su una porta. La probe decodifica la struttura \texttt{sockaddr} all'ingresso della syscall, distinguendo \texttt{AF\_INET} (indirizzi IPv4, con decodifica di porta e indirizzo tramite \texttt{ntop()}) e \texttt{AF\_UNIX} (socket locali, con lettura del percorso). Solo i bind andati a buon fine (valore di ritorno uguale a 0) vengono evidenziati nell'analisi come porte effettivamente aperte.


\noindent\textbf{net\_accept.bt}; traccia \texttt{accept} (id 43) e \texttt{accept4} (id 288), le syscall con cui un processo accetta una connessione in ingresso. Il peer address viene letto dalla struttura \texttt{sockaddr} popolata dal kernel all'uscita della syscall, quando il descrittore della nuova connessione è già disponibile. In fase di analisi, i processi con UID $\geq$ 10000 (UID delle applicazioni Android) che invocano \texttt{accept} vengono segnalati come anomali, poiché su Android i processi applicativi normalmente non agiscono da server.

Una caratteristica comune a tutte e quattro le probe di rete è la pulizia delle mappe temporanee nel blocco \texttt{END}, necessaria per evitare che le strutture dati in memoria nel kernel rimangano allocate al termine del programma bpftrace.

% -------------------------------------------------------
\section{L'orchestratore}

\texttt{monitor.py} è il componente che coordina l'intero ciclo di vita di una sessione di monitoraggio. All'avvio legge la mappa delle probe dal file \texttt{config/probes\_map.json}, verifica che la probe richiesta sia presente nel registro, crea una directory di sessione sotto \texttt{sessions/} con nome composto dal codice della probe e dal timestamp corrente (ad esempio \texttt{syscalls\_2025-01-15\_10-30-00}), e salva immediatamente un file \texttt{session.json} con i metadati della sessione.

\subsection{Gestione del processo bpftrace}

bpftrace viene avviato come sottoprocesso tramite \texttt{subprocess.Popen} con stdout e stderr come pipe separate, in modalità line-buffered. Questa separazione è essenziale: stdout trasporta esclusivamente gli eventi JSON emessi dalle probe tramite \texttt{printf}, mentre stderr contiene i messaggi diagnostici di bpftrace (errori di compilazione, avvisi sui tracepoint, messaggi di avvio). Un thread daemon separato drena continuamente stderr e lo scrive su un file \texttt{stderr.log} all'interno della directory di sessione, garantendo che il thread principale non si blocchi mai in attesa di dati su stderr.

\subsection{Elaborazione degli eventi}

Il thread principale itera riga per riga sullo stdout di bpftrace. Per ogni riga non vuota tenta il parsing JSON; le righe che non sono JSON valido vengono considerate output diagnostico e reindirizzate su \texttt{stderr.log}. Gli eventi JSON validi vengono processati attraverso una pipeline di tre fasi:

\textbf{Normalizzazione}: garantisce che ogni evento rispetti lo schema minimo atteso. Se il campo \texttt{schema\_version} manca, viene aggiunto con il valore definito nella configurazione della probe. Se i campi \texttt{type} o \texttt{event} mancano, vengono aggiunti come fallback dalla configurazione. Se il campo \texttt{data} non è presente o non è un oggetto, viene inizializzato come dizionario vuoto. Viene inoltre aggiunto il campo \texttt{probe\_code} per facilitare la correlazione fuori dalla cartella di sessione.

\textbf{Arricchimento}: attivo solo per gli eventi di tipo \texttt{exec}. Legge \texttt{/proc/<pid>/cmdline} per ottenere la riga di comando completa del processo appena eseguito (i valori sono separati da caratteri null, che vengono sostituiti con spazi). Legge inoltre il link simbolico \texttt{/proc/<pid>/exe} per ottenere il percorso assoluto dell'eseguibile. Questo arricchimento compensa il limite di 15 caratteri del campo \texttt{comm} nel kernel e aggiunge informazioni non accessibili direttamente da bpftrace in un tracepoint.

\textbf{Validazione}: confronta i campi \texttt{type} ed \texttt{event} dell'evento ricevuto con i valori attesi dalla configurazione della probe. Le discrepanze generano avvisi su \texttt{stderr.log} ma non bloccano la scrittura dell'evento, evitando perdita di dati a fronte di probe multi-evento (indicate con \texttt{event: "*"} nella configurazione).

Ogni evento normalizzato, arricchito e validato viene serializzato nuovamente in JSON e scritto su \texttt{events.jsonl} con flush immediato, garantendo che i dati vengano persistiti su disco in tempo reale anche in caso di interruzione improvvisa del processo.

\begin{figure}[h]
    \centering
    \includegraphics[width=0.8\textwidth]{images/monitor-pipeline.png}
    \caption{Pipeline di elaborazione degli eventi in monitor.py}
    \label{fig:monitor-pipeline}
\end{figure}

\subsection{Gestione della sessione e terminazione}

Alla ricezione del segnale di interruzione (Ctrl-C, gestito tramite un blocco \texttt{try/except} su \texttt{KeyboardInterrupt}), \texttt{monitor.py} esegue la pulizia ordinata: termina il processo bpftrace tramite SIGTERM seguito da SIGKILL, aggiorna il file \texttt{session.json} con il timestamp di fine sessione e stampa un riepilogo a terminale. La directory di sessione finale contiene quindi: \texttt{events.jsonl} (gli eventi raccolti), \texttt{stderr.log} (output diagnostico di bpftrace) e \texttt{session.json} (metadati con orari di inizio e fine).

% -------------------------------------------------------
\section{Il launcher interattivo:}

\texttt{launcher.sh} è lo script bash che costituisce il punto di ingresso principale per l'utente. All'avvio esegue una serie di controlli preliminari verificando la presenza di \texttt{python3}, \texttt{bpftrace}, \texttt{jq} e \texttt{timeout} nel PATH, e la disponibilità dei file \texttt{monitor.py} e \texttt{config/probes\_map.json}. In caso di errore viene mostrato un messaggio diagnostico chiaro e lo script termina con codice di uscita 1.

La lista delle probe disponibili viene costruita dinamicamente leggendo le chiavi di \texttt{probes\_map.json} tramite \texttt{jq}, e filtrata tenendo solo le probe il cui file \texttt{.bt} esiste effettivamente su disco. Questa lettura dinamica garantisce che il menu sia sempre sincronizzato con lo stato reale del filesystem, senza necessità di aggiornamenti manuali.

\subsection{Modalità operative}

Lo script supporta tre modalità operative accessibili dal menu principale:

\textbf{Avvio probe}: permette di selezionare una o più probe (con selezione multipla per numero o \texttt{A} per tutte) e un timer in secondi. Le probe vengono eseguite sequenzialmente, ognuna tramite il comando \texttt{timeout} che invia SIGTERM a \texttt{monitor.py} allo scadere del tempo. Durante l'esecuzione viene mostrata una barra di avanzamento ASCII animata che aggiorna stato e percentuale ogni secondo.

\textbf{Generazione report}: mostra una lista delle sessioni disponibili in \texttt{sessions/}, con il numero di eventi raccolti e il timestamp di avvio (letto da \texttt{session.json}). L'utente sceglie la sessione e lo script invoca \texttt{summary.py} per generare il report.

\textbf{Full monitoring}: combina avvio probe e generazione report in sequenza automatica. Dopo la conclusione di ogni probe, il launcher localizza la directory di sessione appena creata estraendo il percorso dall'output di \texttt{monitor.py} tramite \texttt{grep}, con un fallback che cerca la directory più recente in \texttt{sessions/}, e invoca immediatamente \texttt{summary.py} su di essa.

La variabile d'ambiente \texttt{PYTHONUNBUFFERED=1} viene impostata prima di lanciare \texttt{monitor.py} per forzare lo svuotamento immediato del buffer stdout di Python, evitando che gli output vengano persi quando \texttt{timeout} termina il processo prima che i buffer vengano svuotati.

Lo script gestisce anche la modalità non interattiva: se viene passato il flag \texttt{-p} con il percorso di una probe, il menu viene saltato e la probe viene avviata direttamente. Se viene passato anche \texttt{-t}, il timer viene impostato al valore specificato invece di richiedere input dall'utente.

% -------------------------------------------------------
\section{Il modulo di analisi}

\texttt{summary.py} è il componente responsabile dell'analisi post-sessione. Accetta come argomento il percorso di una directory di sessione, legge \texttt{events.jsonl} e \texttt{session.json}, calcola un insieme di metriche aggregate e produce sia un report testuale che un file \texttt{index.json} con tutti i dati analitici in forma strutturata.

\subsection{Analisi generale degli eventi}

La prima fase di analisi calcola statistiche generali: il numero totale di eventi, la loro distribuzione per tipo (\texttt{syscall}, \texttt{process}, \texttt{ipc}, \texttt{network}) e per nome specifico (\texttt{execve}, \texttt{openat}, \texttt{fork}, \texttt{binder\_transaction}, ecc.), e i processi più attivi per volume di eventi generati. Viene inoltre calcolato il tasso di eventi al secondo, sia tramite i metadati di sessione (\texttt{started\_at} e \texttt{stopped\_at} da \texttt{session.json}) sia tramite il campo \texttt{ts\_ns} degli eventi stessi, che permette di identificare il picco di attività al secondo e la finestra temporale più intensa.

\subsection{Analisi delle syscall e della latenza}

Per gli eventi di tipo \texttt{syscall} vengono calcolati conteggi per nome, numero di errori (rilevati tramite valore di ritorno negativo) e tasso di errore percentuale. Se gli eventi includono il campo \texttt{lat\_us} (latenza in microsecondi, presente nelle probe che correlano \texttt{sys\_enter} e \texttt{sys\_exit}), vengono calcolati i percentili p50, p95 e p99 sia globalmente che per nome di syscall e per processo. I 20 eventi più lenti vengono elencati separatamente, e viene prodotta una distribuzione dei codici errno più frequenti, sia globale che per singola syscall.

\subsection{Analisi del Binder IPC}

L'analisi Binder richiede la correlazione di tre tipi di eventi distinti. \texttt{summary.py} costruisce due indici in memoria: uno indicizzato per \texttt{debug\_id} sugli eventi \texttt{binder\_transaction\_alloc\_buf} (per recuperare la dimensione del payload) e uno sugli eventi \texttt{binder\_transaction\_received} (per recuperare il \texttt{comm} del processo destinatario). Per ogni evento \texttt{binder\_transaction} non di tipo reply, viene cercata la corrispondenza nei due indici, costruendo così una visione completa della transazione con mittente, destinatario, dimensione dati e tipo (one-way vs sincrono).

Dal grafo delle transazioni viene costruito un IPC graph che mappa gli archi di comunicazione tra processi, con il numero di chiamate e i byte trasferiti per ogni coppia mittente-destinatario. Questo grafo è il punto di partenza per analisi comportamentali: pattern insoliti di comunicazione, processi che si scambiano volumi anomali di dati, o servizi di sistema contattati da applicazioni non attese.

\begin{figure}[h]
    \centering
    \includegraphics[width=0.9\textwidth]{images/report-binder.png}
    \caption{Esempio di output del report Binder IPC con grafo delle transazioni tra processi.}
    \label{fig:report-binder}
\end{figure}

\subsection{Analisi di rete}

L'analisi di rete aggrega i dati delle quattro probe in un quadro coerente. Vengono costruiti: il port landscape (elenco delle porte effettivamente aperte, filtrato sui bind con \texttt{ret == 0}), la lista dei processi che hanno accettato connessioni in ingresso, il volume di dati inviati e ricevuti per processo, e una timeline al secondo dei byte trasmessi. La timeline viene utilizzata per identificare il picco di invio e il picco di ricezione. I processi con UID $\geq$ 10000 che agiscono da server tramite \texttt{accept} vengono segnalati separatamente come potenzialmente anomali, poiché su Android le applicazioni normalmente non aprono porte in ascolto.

\begin{figure}[h]
    \centering
    \begin{subfigure}{0.48\textwidth}
        \includegraphics[width=\linewidth]{report-rete1.png}
        \caption{Timeline del traffico in uscita}
    \end{subfigure}
    \hfill
    \begin{subfigure}{0.48\textwidth}
        \includegraphics[width=\linewidth]{report-rete2.png}
        \caption{Timeline del traffico in entrata}
    \end{subfigure}
    \caption{Esempio di output del report di analisi di rete}
\end{figure}

\subsection{Albero dei processi e resource map}

A partire dai campi \texttt{pid} e \texttt{ppid} presenti in tutti gli eventi, \texttt{summary.py} ricostruisce l'albero genealogico dei processi osservati durante la sessione. La struttura ad albero viene serializzata sia in forma gerarchica (per la visualizzazione) che in forma flat (per la ricerca rapida). La resource map aggrega invece, per ogni processo, l'insieme dei file aperti (da eventi \texttt{openat}), degli indirizzi di rete contattati (da eventi \texttt{connect}) e dei binari eseguiti (da eventi \texttt{execve}), fornendo una sorta di fingerprint comportamentale per ciascun processo osservato.

\begin{figure}[h]
    \centering
    \begin{subfigure}{0.48\textwidth}
        \includegraphics[width=\linewidth]{albero-processi.png}
        \caption{Albero dei processi}
    \end{subfigure}
    \hfill
    \begin{subfigure}{0.48\textwidth}
        \includegraphics[width=\linewidth]{resource-map.png}
        \caption{Resource map}
    \end{subfigure}
    \caption{Esempio di output del report}
\end{figure}


\subsection{Output}

\texttt{summary.py} produce tre output distinti. Il file \texttt{index.json} nella directory di sessione contiene tutti i dati analitici in forma strutturata e machine-readable, pensato per elaborazioni successive o visualizzazioni esterne. Il report testuale, stampato su stdout e salvato in \texttt{reports/summaries/} con il nome della sessione come identificativo, presenta le stesse informazioni in forma leggibile, con sezioni separate per ogni categoria di analisi e una barra ASCII per la timeline di rete. Il formato del file salvato può essere plain text (\texttt{.txt}) o Markdown (\texttt{.md}), configurabile tramite il flag \texttt{--format}.

% -------------------------------------------------------
\section{Il formato dei dati: JSON Lines}

Tutti gli eventi prodotti dal sistema vengono persistiti~\cite{data_persistence} in formato JSON Lines~\cite{jsonlines} (\texttt{.jsonl}), dove ogni riga è un oggetto JSON indipendente e completo. Questo formato è stato scelto per diverse ragioni pratiche: supporta il parsing incrementale riga per riga senza necessità di caricare l'intero file in memoria, è compatibile con strumenti di analisi standard come \texttt{jq}~\cite{jq_tool}, pandas~\cite{pandas} e Apache Spark~\cite{apache_spark}, e permette lo streaming in tempo reale poiché ogni evento è autonomo e non dipende dal contesto degli eventi precedenti.

Lo schema comune a tutti gli eventi include i campi di intestazione \texttt{ts\_ns}, \texttt{ts}, \texttt{type}, \texttt{event}, \texttt{pid}, \texttt{ppid}, \texttt{tid}, \texttt{uid}, \texttt{comm}, e un oggetto \texttt{data} con i campi specifici della probe. Il campo \texttt{ts\_ns}, espresso in nanosecondi di uptime del kernel, garantisce una risoluzione temporale sufficiente per la correlazione di eventi rapidi come le transazioni Binder. Il campo \texttt{ts}, invece, fornisce un orario leggibile in formato \texttt{HH:MM:SS}, utile per la visualizzazione nei report.


\chapter{Test e Validazione}

In questo capitolo vengono presentati i risultati delle sessioni di monitoraggio effettuate sull'emulatore Android Cuttlefish. Per ciascuno scenario, le sette probe del sistema sono state eseguite in parallelo: la probe delle syscall per il monitoraggio di \texttt{openat}, \texttt{connect} ed \texttt{execve}; la probe del Binder per le transazioni IPC; la probe del ciclo di vita dei processi per la nascita e terminazione; la probe di associazione socket e la probe di accettazione connessioni per il monitoraggio dei socket in ascolto e delle connessioni in ingresso; la probe di invio dati e la probe di ricezione dati per il volume di dati inviati e ricevuti.

Gli scenari analizzati sono tre: l'attivazione della modalità aereo, una sessione di navigazione web con Chrome, e l'apertura dell'applicazione galleria con accesso alla fotocamera.

% -------------------------------------------------------
\section{Attivazione della modalità aereo}
\label{sec:aereo}

L'applicazione monitorata è \texttt{com.android.settings} (uid 1000), il pannello di configurazione di sistema di Android. Lo scenario consiste nell'attivare e poi disattivare la modalità aereo tramite il menu delle impostazioni rapide. Il riferimento temporale è l'avvio della probe del Binder, con le altre sei probe avviate contestualmente. I due eventi chiave sono l'attivazione della modalità aereo e la sua disattivazione.

\begin{figure}[h]
    \centering
    \includegraphics[width=\textwidth]{images/timeline-aereo.png}
    \caption{Sequenza delle azioni nella sessione di attivazione e disattivazione della modalità aereo.}
    \label{fig:timeline-aereo}
\end{figure}

\subsection{Fase preliminare}

Nei 100 secondi precedenti all'attivazione, il sistema opera nella schermata home con il WiFi attivo. La probe del Binder registra un'attività di fondo dominata dal canale chiamato da \texttt{surfaceflinger}, con 3.330 chiamate per circa 1.4~MB: è il compositor grafico~\cite{surfaceflinger_docs} che aggiorna continuamente la schermata.

\subsection{Attivazione della modalità aereo}

All'attivazione della modalità aereo si produce un effetto immediato e netto sullo stack di rete. La probe di associazione socket non registra alcun evento tra durante l'intera finestra che copre il periodo in modalità aereo e la successiva riattivazione non genera alcun \texttt{bind} da processi di rete. È la conferma diretta che lo stack WiFi è stato completamente fermato.

Il percorso del comando attraverso il sistema è ricostruibile dalle probe. L'utente interagisce con \texttt{com.android.settings}, che invia via Binder~\cite{binder_ipc} il \textit{broadcast} di sistema \texttt{ACTION\_AIRPLANE\_MODE\_CHANGED} al \texttt{ConnectivityService} all'interno del \textit{system server}. È quest'ultimo, non l'applicazione impostazioni, a coordinare la sequenza di disattivazione: WiFi e Bluetooth vengono fermati attraverso chiamate ai rispettivi HAL via Binder. Il processo \texttt{ndroid.settings} accumula 4.169 transazioni Binder nell'intera sessione, ma nessuna di esse produce direttamente traffico di rete: tutto passa attraverso il \textit{system server}.

\subsection{Disattivazione e riconnessione}

Alla disattivazione della modalità aereo viene avviata una sequenza di riconnessione osservabile passo per passo attraverso le probe. Il primo segnale misurabile arriva con un po' di ritardo rispetto al toggle: vengono registrate le prime operazioni \texttt{bind} dopo la pausa. Successivamente \texttt{IpClient.wlan0} apre sei nuovi socket per avviare la negoziazione DHCP, e \texttt{netd} crea un thread per la risoluzione dei nomi. Il ritardo riflette il tempo necessario al \textit{firmware} dell'interfaccia WiFi virtuale di Cuttlefish per completare la riassociazione alla rete prima che lo stack di rete superiore possa riattivarsi.

La probe di ricezione dati mostra che \texttt{IpClient.wlan0} riceve complessivamente circa 229~KB nell'intera sessione: è il traffico DHCP~\cite{dhcp_rfc} corrispondente all'acquisizione dell'indirizzo IP dopo la riconnessione.

Il picco di Binder dell'intera sessione (4.099 eventi/s) si registra due secondi prima del primo bind di rete. Esso è prodotto dai \textit{broadcast} di \texttt{ConnectivityService} che notificano tutti i processi del sistema del cambiamento di stato della rete, prima ancora che il collegamento fisico sia ristabilito. Anche questa sequenza, il Binder che esplode prima che i socket di rete si riaprino, è osservabile solo grazie al monitoraggio parallelo delle sette probe.

\subsection{Osservazioni sull'architettura osservata}
\label{subsec:aereo-arch}

Lo scenario della modalità aereo si distingue rispetto ai prossimi due per la natura dell'evento monitorato: non si tratta dell'avvio di un'applicazione ma di un cambiamento di stato del sistema che coinvolge l'infrastruttura di rete.

Il segnale più caratteristico è la finestra di silenzio nelle probe di rete. La probe di associazione socket non registra eventi, e la probe di accettazione connessioni mostra un gap analogo. Questa assenza prolungata di attività, in un sistema che normalmente produce bind e accept continuamente, è l'indicatore diretto della disattivazione dello stack WiFi.

Il traffico inviato via socket (2.8~MB totali) è dominato non dal WiFi ma dalle comunicazioni interne dell'hardware audio e grafico (\texttt{android.hardwar}: 1.8~MB, \texttt{app} SurfaceControl: 709~KB). Il contributo diretto di \texttt{wificond} è di soli 4.9~KB, segno che il controllo della radio WiFi avviene quasi interamente via Binder e non via socket.

\begin{table}[h]
\centering
\begin{tabular}{lr}
\toprule
\textbf{Categoria} & \textbf{Volume totale} \\
\midrule
Transazioni Binder & 35.058 \\
Traffico IPC & 11.3 MB \\
Dati inviati via socket & 2.8 MB \\
Dati ricevuti via socket & 4.8 MB \\
Bind durante modalità aereo & 0 \\
\bottomrule
\end{tabular}
\caption{Riepilogo del traffico complessivo della sessione modalità aereo: il silenzio nelle probe di rete durante la finestra t$\,\approx\,$92s--147s è il segnale diretto della disattivazione dello stack WiFi.}
\end{table}

% -------------------------------------------------------
\section{Navigazione web con Chrome}
\label{sec:chrome}

L'applicazione monitorata è \texttt{org.chromium.webview\_shell} (uid 10064), un'applicazione di test che espone direttamente il motore di rendering WebView di Chromium. Lo scenario consiste nella navigazione web a partire dalla schermata home dell'emulatore: viene prima effettuata una ricerca tramite il widget di ricerca presente nella home, poi si naviga sul sito google.com e si esegue una ricerca di notizie. Le probe vengono avviate prima di qualsiasi interazione con il browser, per poter documentare il comportamento di base del sistema e distinguerlo dall'attività generata dall'utente.

La sessione si articola in quattro fasi, identificate in base alle azioni effettuate:

\begin{figure}[h]
    \centering
    \includegraphics[width=\textwidth]{images/timeline-rete.png}
    \caption{Sequenza delle azioni nella sessione di navigazione web.}
    \label{fig:timeline-chrome}
\end{figure}

\subsection{Fase preliminare}

Nei primi 36 secondi l'emulatore si trova nella schermata home e non viene effettuata alcuna interazione. Questa fase serve a documentare il comportamento di base del sistema: Android non è silenzioso nemmeno in \textit{idle}, e senza questa \textit{baseline} sarebbe impossibile distinguere il traffico di fondo dall'attività riconducibile alle azioni utente.

La probe di ricezione dati mostra traffico già dai primi secondi, generato da \texttt{logcat} e da altri servizi di sistema che comunicano su socket locali in modo continuativo. Più interessante è un picco di invio che compare intorno a t$\,\approx\,$27s, con il sistema che trasmette prima 6.4~KB poi 28.9~KB in due secondi consecutivi. Il traffico è dominato da \texttt{InputDispatcher} (18.4~KB a t$\,\approx\,$28s) e dal processo \texttt{app} (SurfaceControl), con un contributo minore di \texttt{android.ui}; non ha nulla a che fare con il browser, che non è ancora avviato. Si tratta di comunicazioni interne al sottosistema grafico e di input, un esempio di quanto il sistema operativo produca traffico di rete per ragioni completamente indipendenti dall'utente.

\begin{figure}[h]
    \centering
    \includegraphics[width=\textwidth]{images/picco-invii.png}
    \caption{Picco di invio nella fase preliminare, generato dal compositor grafico prima di qualsiasi interazione utente.}
    \label{fig:picco-invii}
\end{figure}

\subsection{Invio della query}

Il click sul widget di ricerca non produce eventi netti nelle probe di rete. L'attività misurabile parte qualche secondo dopo, quando la probe delle syscall intercetta una \texttt{connect} da parte dell'applicazione Google Search che gestisce il widget nella schermata home. Contestualmente, la probe di accettazione connessioni mostra che \texttt{netd} crea un thread che  nasce per gestire la sua richiesta DNS. È quindi l'app di ricerca Google, e non il browser, ad avviare la risoluzione del nome mentre l'utente digita la query.

La prima attività di rete del processo Chrome compare circa nove secondi dopo il click, quando le probe di invio e ricezione dati registrano i primi eventi dal thread \texttt{RenderThread} (PID 8373). I thread \texttt{m.webview\_shell} e \texttt{NetworkService} compaiono nei secondi successivi. Il fatto che thread con nomi diversi condividano lo stesso PID è un effetto del limite di 15 caratteri del campo \texttt{comm} nel kernel Linux~\cite{linux_task_struct}, che riporta il nome del singolo thread piuttosto che quello del processo.

Successivamente, la probe del Binder registra i primi eventi associati a \url{CrRenderer\_Main} (PID 8406),
il renderer di Chromium. La comunicazione tra \texttt{NetworkService}
e il renderer avviene via Binder, non via socket: nell'arco dell'intera sessione vengono registrate 1031 transazioni tra i 
due per un totale di circa 339~KB. Ogni risposta HTTP ricevuta dalla rete viene consegnata al renderer attraverso questo 
canale IPC prima di essere visualizzata.

\subsection{Caricamento di google.com}

Il click su google.com produce in un secondo un \textit{burst} di otto operazioni \texttt{bind} da parte di \texttt{NetworkService}, che indicano l'apertura di nuovi socket per il caricamento parallelo delle risorse della pagina, coerente con il \textit{multiplexing} HTTP/2~\cite{http2_rfc} utilizzato da Chrome.

Il picco di connessioni in ingresso arriva con un ritardo di circa dodici secondi rispetto al click. Il ritardo è compatibile con il tempo di caricamento del contenuto dinamico della homepage di Google, che include numerose richieste secondarie per immagini, script e dati personalizzati.

\subsection{Ricerca ``google news''}

La ricerca di ``google news'' produce la risposta più netta dell'intera sessione. La probe delle syscall registra il picco assoluto di 608 eventi al secondo, contro una media di circa 187 per l'intera sessione.

Il picco di invio dati segue con un ritardo di circa dodici secondi: la probe di invio dati registra 79.4~KB inviati in un secondo, il valore massimo dell'intera sessione. Il ritardo riflette il tempo che il browser impiega a ricevere i risultati, elaborarli tramite il renderer e inviare le richieste secondarie per i contenuti collegati.

\begin{figure}[h]
    \centering
    \includegraphics[width=\textwidth]{images/send-peak.png}
    \caption{Picco di 79.4~KB inviati in un secondo, valore massimo dell'intera sessione.}
    \label{fig:send-peak}
\end{figure}

\subsection{Osservazioni sull'architettura osservata}
\label{subsec:chrome-arch}

Il monitoraggio parallelo delle sette probe ha permesso di osservare una catena di eventi che nessuna singola probe avrebbe potuto documentare per intero.

Il renderer \texttt{CrRendererMain} (PID 8406) non ha accesso diretto alla rete. Tutta la comunicazione tra il codice di rendering e lo stack di rete passa attraverso Binder: la catena osservata è \texttt{m.webview\_shell} $\rightarrow$ Binder $\rightarrow$ \texttt{NetworkService} $\rightarrow$ socket. Questo isolamento è una scelta progettuale di Chromium~\cite{chromium_multiprocess} per separare il processo di rendering, potenzialmente esposto a contenuti non fidati, dal processo che gestisce le connessioni di rete.

Complessivamente, la sessione ha prodotto 42\,343 transazioni Binder per un totale di circa 17~MB trasferiti via IPC, a fronte di 1.55~MB inviati e 2.37~MB ricevuti via socket. Il volume di traffico IPC interno supera quindi di circa cinque volte il traffico di rete esterno: anche una semplice sessione di navigazione web genera internamente un volume di comunicazione inter-processo di gran lunga superiore a quello scambiato con la rete.

\begin{table}[h]
\centering
\begin{tabular}{lr}
\toprule
\textbf{Categoria} & \textbf{Volume totale} \\
\midrule
Transazioni Binder & 42.343 \\
Traffico IPC & 17 MB \\
Dati inviati via socket & 1.55 MB \\
Dati ricevuti via socket & 2.37 MB \\
\bottomrule
\end{tabular}
\caption{Riepilogo del traffico complessivo della sessione Chrome: 42\,343 transazioni Binder e confronto tra volume IPC e traffico di rete.}
\end{table}

% -------------------------------------------------------
\section{Apertura della galleria e scatto fotografico}
\label{sec:galleria}

L'applicazione monitorata è l'applicazione predefinita per la gestione delle immagini su Android (uid 10066). Lo scenario combina due applicazioni distinte: nella prima parte l'utente naviga la galleria fotografica, nella seconda passa alla fotocamera e scatta un'immagine. Il riferimento temporale è l'avvio della probe del Binder, con le altre sei probe avviate contestualmente nei secondi precedenti. Gli eventi si distribuiscono in tre momenti precisi: apertura della galleria, attivazione della fotocamera e scatto della foto.

\begin{figure}[h]
    \centering
    \includegraphics[width=\textwidth]{images/timeline-galleria.png}
    \caption{Sequenza delle azioni nella sessione di apertura galleria e fotocamera.}
    \label{fig:timeline-galleria}
\end{figure}

\subsection{Apertura della galleria}

All'apertura dell'applicazione galleria, la probe del ciclo di vita dei processi registra la creazione del processo \texttt{droid.gallery3d} (PID 14805) come figlio del \textit{system server} di Android. L'associazione all'uid 10066 è confermata dai percorsi dei file \texttt{cgroup} che il processo \texttt{ActivityManager} monitora per la gestione della memoria dell'applicazione.

La probe delle syscall rileva le prime operazioni di accesso allo storage: un thread dell'applicazione apre la radice dello spazio di archiviazione condiviso. Questo accesso corrisponde alla scansione iniziale del catalogo di immagini che la galleria esegue all'avvio per costruire la lista dei contenuti disponibili.

Sul fronte del Binder, l'avvio della galleria produce un'attività crescente che culmina nel picco assoluto della sessione: la probe registra 4\,723 transazioni in un singolo secondo, circa cinque volte la media della sessione (941 eventi/s). Questo picco corrisponde al momento in cui la galleria carica le anteprime delle immagini, operazione che richiede un'intensa sequenza di comunicazioni con \texttt{SurfaceFlinger}~\cite{surfaceflinger_docs} per il rendering di ciascun thumbnail.

\begin{figure}[h]
    \centering
    \includegraphics[width=\textwidth]{images/binder-galleria-peak.png}
    \caption{Picco di 4\,723 transazioni Binder, durante il caricamento delle anteprime da parte della galleria.}
    \label{fig:binder-galleria-peak}
\end{figure}

\subsection{Attivazione della fotocamera}

Il passaggio alla fotocamera produce il cambiamento più caratteristico dell'intera sessione nel profilo delle comunicazioni Binder. Il processo \texttt{android.camera2} (PID 1913) diventa attivo e due thread legati all'hardware della fotocamera emergono tra i processi più trafficati: \texttt{Cam0\_ResultDisp} con 6\,285 eventi e \texttt{C3Dev-0-ReqQueu} con 4\,280 eventi, il secondo e il quarto posto nell'intera sessione.

Il grafo IPC mostra che \texttt{Cam0\_ResultDisp} comunica principalmente con il canale che corrisponde alla pipeline di elaborazione dei frame della fotocamera: il thread che riceve i risultati dal sensore li trasmette al \texttt{Camera Hardware Abstraction Layer}~\cite{camera_hal} attraverso Binder, che li consegna al processo di visualizzazione per l'anteprima in tempo reale. L'attivazione di questo canale è assente in tutti gli altri scenari analizzati e rappresenta quindi un segnale affidabile dell'accesso all'hardware della fotocamera.

In parallelo, \texttt{SurfaceFlinger} mantiene il canale più trafficato dell'intera sessione: 4\,481 chiamate verso il \texttt{binder} per un totale di circa 2.1~MB. Questo traffico riflette la composizione continua dei frame dell'anteprima, che richiede una frequenza di comunicazione elevata tra il compositor grafico e il driver del display.

\begin{figure}[h]
    \centering
    \includegraphics[width=\textwidth]{images/binder-camera-ipc.png}
    \caption{Grafo delle comunicazioni Binder durante l'anteprima della fotocamera: il canale \texttt{Cam0\_ResultDisp} verso il Camera HAL è il segnale caratteristico dell'accesso all'hardware fotografico.}
    \label{fig:binder-camera-ipc}
\end{figure}

\subsection{Scatto della foto}

Lo scatto della fotografia  lascia una traccia diretta nella resource map prodotta dalla probe delle syscall. Un thread del processo \texttt{droid.gallery3d} apre il file \texttt{IMG\_20260304\_113342.jpg}: si tratta dell'immagine appena acquisita, letta dall'applicazione per aggiornare il catalogo con il nuovo scatto. Lo stesso thread accede anche a \texttt{/share\_history.xml}, il file che l'applicazione aggiorna sistematicamente a ogni modifica del catalogo.

La probe di ricezione dati registra circa 215~KB ricevuti dal processo \texttt{droid.gallery3d}  su socket locali. Il ritardo di circa tredici secondi rispetto allo scatto è compatibile con il tempo necessario al sistema di gestione dei media di Android per processare il nuovo file, generare l'anteprima e notificare l'applicazione dell'avvenuto aggiornamento del catalogo.

Nessuna comunicazione di rete esterna è associata allo scatto: la probe di invio dati registra per \texttt{droid.gallery3d} un totale di soli circa 4~KB durante l'intera sessione, esclusivamente su socket locali verso servizi di sistema.

\subsection{Osservazioni sull'architettura osservata}
\label{subsec:galleria-arch}

Lo scenario galleria e fotocamera si distingue dalla navigazione web per l'assenza totale di traffico di rete esterno. La probe di associazione socket non registra alcun evento per l'intera sessione, e quella di accettazione connessioni rileva attività solo da processi di sistema come \texttt{init} e \texttt{netd}, mai dalle applicazioni utente. È la conferma che galleria e fotocamera operano esclusivamente su risorse locali.

Nonostante questa chiusura alla rete, il volume di comunicazione IPC è considerevole: la sessione ha prodotto 38\,824 transazioni Binder per un totale di circa 17.2~MB trasferiti via IPC, a fronte di appena 4~KB inviati via socket. Anche operazioni che l'utente percepisce come completamente locali generano internamente un flusso continuo di messaggi tra decine di componenti di sistema.

Il dato più interessante riguarda l'identificazione dell'accesso alla fotocamera. L'attivazione del thread \texttt{Cam0\_ResultDisp} e del canale Binder verso il Camera HAL è un indicatore preciso e stabile: non compare durante l'uso della galleria, appare non appena la fotocamera viene avviata, e persiste fino alla chiusura dell'applicazione. Questo tipo di segnatura IPC potrebbe essere utilizzato per rilevare l'accesso alla fotocamera da parte di un processo, indipendentemente da qualsiasi indicatore a livello applicativo.

\begin{table}[h]
\centering
\begin{tabular}{lr}
\toprule
\textbf{Categoria} & \textbf{Volume totale} \\
\midrule
Transazioni Binder & 38\,824 \\
Traffico IPC & 17.2 MB \\
Dati inviati via socket & 4 KB \\
Dati ricevuti via socket & 215 KB \\
Bind aperti & 0 \\
\bottomrule
\end{tabular}
\caption{Riepilogo del traffico complessivo della sessione galleria e fotocamera: 38\,824 transazioni Binder e traffico di rete trascurabile.}
\end{table}

\chapter{Limiti e possibili miglioramenti}

Il sistema sviluppato raggiunge gli obiettivi dichiarati nel capitolo introduttivo: permette di eseguire bpftrace su Android, raccogliere eventi da sette probe parallele e analizzare il comportamento delle applicazioni a livello di kernel. Le sperimentazioni descritte nel capitolo precedente dimostrano che i dati prodotti sono sufficientemente precisi e ricchi da permettere ricostruzioni dettagliate di scenari reali. Rimangono tuttavia limiti concreti che ne condizionano l'applicabilità, e margini di miglioramento che potrebbero ampliarne l'utilità in contesti diversi da quello in cui è stato sviluppato.

% -------------------------------------------------------
\section{Limiti del sistema attuale}

\subsection{Dipendenza dall'ambiente emulato}

Il vincolo più significativo è che l'intero sistema è progettato per funzionare su Android Cuttlefish in esecuzione su Ubuntu, con una distribuzione Debian installata all'interno dell'emulatore. Questo ambiente è riproducibile e controllato, ma non corrisponde a un dispositivo reale. Su un dispositivo Android di produzione, l'esecuzione di bpftrace richiederebbe accesso root, che a sua volta implica il rooting del dispositivo o la sostituzione della ROM con una versione personalizzata. Entrambe le opzioni sono impraticabili nella maggior parte dei contesti operativi reali, come documentato nella Tabella~\ref{fig:confronto-ambienti} del capitolo introduttivo.

Anche laddove l'accesso root fosse disponibile, l'ambiente Debian necessario all'esecuzione di bpftrace non sarebbe presente, e installarlo su un dispositivo fisico richiederebbe una procedura significativamente più complessa di quella seguita per Cuttlefish. Il sistema è pertanto utile per analisi in ambienti di sviluppo e test, ma non è distribuibile su flotte di dispositivi reali nella forma attuale.

\subsection{Analisi esclusivamente post-sessione}

Il sistema attuale è progettato per raccogliere dati durante una sessione e analizzarli al termine. Non è possibile ricevere notifiche in tempo reale su eventi rilevanti, né interrompere automaticamente la sessione in risposta a comportamenti anomali osservati. Questo limita l'utilità del sistema in scenari di monitoraggio continuo o di risposta agli incidenti, dove la tempestività dell'informazione è importante.

\subsection{Identificazione dei processi}

Il limite di 15 caratteri del campo \texttt{comm} nel kernel Linux rende difficile l'identificazione univoca dei processi nelle probe. Come documentato nelle analisi del capitolo precedente, nomi come \texttt{ndroid.settings}, \texttt{ndroid.systemui} o \texttt{droid.gallery3d} sono troncamenti che richiedono conoscenza del sistema per essere ricondotti all'applicazione corretta. Il modulo di orchestrazione mitiga parzialmente questo problema leggendo \texttt{/proc/<pid>/cmdline} all'evento \texttt{exec}, ma questa arricchimento non è sempre disponibile per i processi già in esecuzione al momento dell'avvio delle probe.

% -------------------------------------------------------
\section{Direzioni di sviluppo}

\subsection{Estensione a dispositivi reali}

La direzione più impattante per ampliare l'applicabilità del sistema è rendere le probe eseguibili su dispositivi Android reali. Una via percorribile senza rooting del dispositivo è l'utilizzo di Android Debug Bridge (\texttt{adb}) in modalità \textit{shell} con i privilegi elevati disponibili nelle versioni di debug del sistema operativo, distribuite per scopi di sviluppo. In questo scenario le probe potrebbero essere caricate dall'esterno tramite \texttt{adb shell} senza richiedere modifiche permanenti al dispositivo.

Un'altra direzione è l'integrazione con ambienti di test aziendali che già operano con ROM personalizzate o dispositivi con accesso root abilitato a scopo di analisi. In questi contesti il sistema potrebbe essere distribuito come strumento di analisi della sicurezza senza le limitazioni tipiche dei dispositivi di produzione.

\subsection{Analisi in tempo reale e sistema di allerta}

L'architettura a \textit{pipeline} descritta nel capitolo di implementazione si presta a un'estensione verso l'analisi in tempo reale. Invece di scrivere gli eventi su file \texttt{.jsonl} e analizzarli al termine della sessione, il modulo di orchestrazione potrebbe esporre un flusso di eventi a un componente di analisi continua. Quest'ultimo potrebbe applicare regole di rilevamento e generare notifiche al verificarsi di condizioni predefinite, come l'apertura di connessioni verso indirizzi IP non visti in precedenza, l'esecuzione di binari fuori dalle directory di sistema, o l'attivazione di canali Binder verso il Camera HAL da parte di un'applicazione non autorizzata.

La soglia di frequenza di campionamento delle probe attuali è sufficiente per questo tipo di analisi: come mostrato nel capitolo precedente, i picchi di attività rilevanti sono ben distinguibili dal traffico di fondo già nell'arco di pochi secondi.

\subsection{Firme comportamentali e rilevamento di anomalie}

I dati raccolti nelle sessioni di test evidenziano che le applicazioni Android producono profili di attività caratteristici e riproducibili: la galleria genera un picco Binder specifico al caricamento delle anteprime, la fotocamera attiva thread HAL identificabili, il browser produce sequenze di \texttt{bind} correlate alle richieste HTTP/2. Questi profili potrebbero essere formalizzati come firme comportamentali da confrontare con l'attività osservata in sessioni successive.

Un sistema di rilevamento basato su firme potrebbe segnalare, ad esempio, un'applicazione che attiva il canale Binder verso il Camera HAL pur non essendo un'applicazione fotografica, o un processo che apre socket verso l'esterno in assenza di interazione utente. Questo tipo di rilevamento non richiederebbe tecniche di apprendimento automatico nella fase iniziale: le firme osservate nelle sperimentazioni sono già sufficientemente precise da definire regole basate su soglie e pattern.

\subsection{Estensione delle probe}

Le sette probe attuali coprono syscall di rete, Binder IPC e ciclo di vita dei processi. Aree di sistema attualmente non monitorate che potrebbero offrire informazioni complementari includono:

\begin{itemize}
    \item la gestione della memoria, tramite il tracciamento di \texttt{mmap} e \texttt{brk}, che permetterebbe di correlare picchi di allocazione con l'attività delle applicazioni;
    \item la schedulazione dei processi, tramite i tracepoint \texttt{sched:sched\_switch} e \texttt{sched:sched\_wakeup}, utile per misurare la contesa di CPU tra componenti di sistema;
    \item le operazioni di lettura e scrittura su file, tramite \texttt{read} e \texttt{write}, per quantificare il volume di I/O su storage e non solo gli accessi rilevati da \texttt{openat}.
\end{itemize}

L'architettura modulare del sistema facilita l'aggiunta di nuove probe senza modificare i componenti esistenti: è sufficiente implementare il nuovo script bpftrace, registrarlo nel file di configurazione condiviso, e il modulo di orchestrazione e l'interfaccia utente lo renderanno automaticamente disponibile.

\subsection{Interfaccia di visualizzazione}

Il formato attuale dell'output è un report testuale generato dal modulo di analisi. Per sessioni lunghe o con molte probe attive simultaneamente, la lettura dei report diventa onerosa. Un'interfaccia di visualizzazione basata su linea temporale, anche semplice come un file HTML generato staticamente, permetterebbe di sovrapporre gli eventi delle diverse probe su un asse temporale comune e identificare visivamente le correlazioni tra probe. La disponibilità dei dati in formato JSON strutturato, già prodotto dal modulo di analisi, renderebbe questo sviluppo relativamente agevole senza modifiche al nucleo del sistema.

\chapter{Conclusioni}

Questa tesi ha affrontato il problema dell'utilizzo di bpftrace per il tracing di applicazioni Android, partendo dalla constatazione che gli strumenti di osservabilità tipici dei sistemi Linux non sono direttamente disponibili su Android senza adattamenti significativi. La soluzione proposta, basata sull'esecuzione di bpftrace all'interno di una distribuzione Debian installata nell'emulatore Cuttlefish, ha dimostrato che è possibile raccogliere dati di sistema a livello kernel su Android senza modificare il sistema operativo e senza richiedere accesso hardware fisico.

Il sistema sviluppato comprende sette probe bpftrace parallele, un modulo di orchestrazione in Python che coordina la raccolta degli eventi, un modulo di analisi post-sessione e un'interfaccia utente bash. Le probe coprono il ciclo di vita dei processi, le syscall di accesso a file e rete, le transazioni IPC Binder e il volume di dati scambiati via socket. L'architettura a pipeline consente di attivare solo il sottoinsieme di probe necessario per lo scenario di interesse, limitando il volume di dati prodotti e l'overhead sul sistema monitorato.

Le sperimentazioni condotte su tre scenari distinti hanno mostrato che i dati raccolti dal sistema sono sufficientemente precisi da permettere ricostruzioni dettagliate del comportamento delle applicazioni a livello di sistema operativo. Nella sessione di navigazione web con Chrome è stato possibile osservare la catena completa che collega l'interazione dell'utente al traffico di rete: la risoluzione DNS avviata dall'applicazione di ricerca Google, il passaggio dei dati HTTP dal \textit{NetworkService} al renderer Chromium via Binder, e i picchi di attività correlati a ciascuna fase della navigazione. Il volume di traffico IPC interno, circa 17~MB, ha superato di cinque volte il traffico di rete esterno, un dato che da solo non sarebbe emerso dall'analisi di una singola probe.

Nella sessione galleria e fotocamera, il sistema ha permesso di identificare la firma Binder specifica dell'accesso all'hardware fotografico: l'attivazione del thread \texttt{Cam0\_ResultDisp} e del canale verso il Camera HAL è risultata assente durante la sola navigazione della galleria e presente non appena la fotocamera è stata avviata. Questa firma è stabile e riproducibile, e non dipende da indicatori a livello applicativo.

Nella sessione della modalità aereo, la finestra di silenzio nelle probe di associazione socket e accettazione connessioni ha documentato con precisione la disattivazione dello stack WiFi, e la sequenza di riconnessione successiva è risultata ricostruibile passo per passo attraverso i bind di \texttt{wpa\_supplicant}, la negoziazione DHCP di \texttt{IpClient.wlan0} e la creazione dei thread DNS di \texttt{netd}. Il picco Binder che precede di due secondi la ricomparsa dei primi bind di rete ha evidenziato come il \textit{system server} propaghi i cambiamenti di stato della rete a tutti i processi del sistema prima che il collegamento fisico sia effettivamente ristabilito.

In tutti e tre gli scenari, le correlazioni più significative sono emerse dal confronto tra probe diverse, non dall'analisi di una singola fonte di dati. Questo risultato conferma la premessa del lavoro: un sistema di monitoraggio parallelo e multidimensionale permette di osservare il comportamento di Android in modo qualitatively diverso rispetto a strumenti che operano su una sola categoria di eventi.

I limiti del sistema, discussi nel capitolo precedente, riguardano principalmente la dipendenza dall'ambiente emulato e l'assenza di analisi in tempo reale. Entrambi sono limiti architetturali che non invalidano il contributo della tesi, ma definiscono il confine dell'applicabilità attuale: il sistema è uno strumento di analisi e ricerca, non un agente di monitoraggio distribuibile su dispositivi di produzione nella forma attuale.

Il contributo principale di questo lavoro è dimostrare che bpftrace, con l'infrastruttura appropriata, può essere utilizzato per ottenere una visibilità profonda sul comportamento di Android a livello kernel. I dati raccolti hanno permesso di osservare meccanismi interni del sistema operativo, come l'architettura di isolamento del renderer Chromium, la pipeline di elaborazione dei frame della fotocamera e la sequenza di riconnessione WiFi, che non sarebbero accessibili attraverso le API applicative o gli strumenti di profilazione standard. Questo apre prospettive concrete per l'utilizzo di tecniche eBPF nell'analisi della sicurezza delle applicazioni Android, nella verifica del comportamento di componenti di sistema e nel rilevamento di anomalie basato su firme comportamentali a livello kernel.

\input{chapters/chap07.tex}

\nocite{*}
\printbibliography[heading=bibintoc]

\end{document}
